\chapter{Discusión y conclusiones}
\label{cap:discusion}

\section{Discusión de los resultados}

\subsection{Capacidad de arbitraje y seguimiento de precios (OE2, H1)}

En una piscina aislada de producto constante, el resultado del arbitraje se deriva directamente de la ecuación de intercambio: la capacidad $c$ mueve el precio logarítmico hasta en $2\log(1+c)$ por tick hasta llegar a la banda de no arbitraje. Para el choque logarítmico ensayado de $0.20$, la capacidad baja $c=0.005$ requiere 20 ticks ideales para corregir el exceso fuera de la banda. Por tanto, la reducción simulada del área de exceso confirma la implementación bajo un flujo de órdenes interactivo, pero su dirección no constituye un descubrimiento independiente. Su contenido cuantitativo más útil es el patrón de recuperación en horizonte finito y la diferencia entre el referente aislado y el calendario con ruido.

\subsection{Retiro de liquidez y sensibilidad H2 (OE3)}

La regla fija de retiro se activa en todas las corridas E3 porque cada choque ensayado supera su umbral; es una respuesta única y no representa conductas heterogéneas de proveedores. La repetición de las semillas existentes aclara el mecanismo en la celda poco profunda. En E1-alta estático, las 143 salidas posteriores observadas son episodios únicos que coinciden con el tick configurado de pánico, y la corrida pareada sin pánico no presenta salidas. En E3, el 98.2--99.6\,\% de las trayectorias recuperadas vuelve a salir, normalmente en episodios múltiples, y el patrón se mantiene sin pánico. Un control pareado que fija el límite de entrada del arbitraje en su valor basado en las reservas iniciales reduce el contraste E3 menos estático del área de exceso posrecuperación de 0.122--0.124 a 0.0036--0.0097. Esto indica que la capacidad ligada a reservas explica buena parte del contraste en el diseño seleccionado. El control es posterior al análisis inicial, reutiliza semillas, condiciona el área a la recuperación y cubre una profundidad; no aísla un efecto causal universal del retiro o la profundidad. H2 es una interacción condicional entre la regla de umbral, el pánico, el tamaño de órdenes, las reservas y la capacidad.

\subsection{Mapas reducidos y alcance predictivo (OE4, H3)}

Esto responde la cláusula del mapa como descripción predictiva limitada, no como clasificación de regímenes: el mapa de dos estados mejora la persistencia en algunas condiciones reservadas, pero no supera al AR(1) escalar, por lo que su segundo estado no demuestra valor predictivo incremental a un paso. La mejora media de alrededor de 1.7\,\% del error de reproducción del mapa sensible a la banda se concentra en una condición de baja capacidad/profundidad y depende de la trayectoria futura realizada del CEX. Es un resultado condicional de reproducción, no un pronóstico en tiempo real.

\section{Conclusiones por objetivo}

\begin{description}
  \item[OE1.] El simulador pasó las invariantes de implementación y sus resultados registrados son reproducibles. Esto establece consistencia con las reglas codificadas.
  \item[OE2.] La corrección ideal según la capacidad se deriva exactamente. La simulación muestra las consecuencias en horizonte finito de añadir el ruido, la tendencia y el calendario especificados; no ofrece una prueba independiente de que el arbitraje mueva el precio hacia la referencia.
  \item[OE3.] La ablación del pánico distingue las salidas E1 únicas asociadas al pánico de las salidas E3 recurrentes que persisten sin pánico. En la profundidad baja ensayada, fijar el límite de arbitraje en las reservas iniciales elimina la mayor parte del contraste posrecuperación de E3. Estos resultados identifican un mecanismo dentro de este simulador, condicionado a la regla, el calendario de órdenes y la selección por recuperación; no establecen un efecto general del retiro.
  \item[OE4.] El mapa de dos estados no mejora el AR(1) escalar. El mapa sensible a la banda tiene una ventaja de reproducción pequeña y localizada que utiliza la referencia futura realizada; no hay evidencia de predicción operativa.
\end{description}

\section{Conclusión general}

Dentro de la grilla especificada, el cálculo de piscina aislada da la corrección ideal por evento, mientras que la trayectoria simulada posperturbación también depende de la profundidad de reservas, el tamaño de las órdenes, las reglas de agentes y el orden de eventos. La pregunta revisada se refiere a estas trayectorias finitas de error y no a regímenes de mercado. Para ese objetivo, el mapa de dos estados no supera al AR(1) escalar, y la pequeña ganancia de reproducción del mapa sensible a la banda es localizada y usa la trayectoria futura realizada del CEX. Por tanto, la evidencia responde la pregunta revisada dentro de este simulador, pero no establece un efecto de mercado real, un pronóstico operativo ni un clasificador de regímenes.

El resultado de persistencia de la recuperación se concentra en E3. En sus nueve celdas de profundidad y choque ponderadas por igual, el 31.2\,\% de las trayectorias que logró una recuperación sostenida salió después de la banda; con factor de profundidad 0.5, la proporción condicional fue 98.8\,\%, mientras que con factores 1 y 2 fue cero. El 31.2\,\% es un resumen ponderado del diseño del simulador, no una frecuencia de mercado estimada. En las celdas poco profundas, la repetición distingue dos patrones: las salidas posteriores de E1 son episodios únicos asociados al pánico, mientras que las de E3 se repiten incluso sin pánico. El control de capacidad fija atenúa sustancialmente el contraste de E3. En futuras reproducciones de eventos se deberían registrar por separado el primer ingreso a banda, el retorno sostenido de dos ticks, las salidas posteriores, los eventos de pánico/órdenes, las reservas activas y la capacidad de arbitraje disponible; son observables candidatos, no requisitos validados para monitoreo en vivo.

El estimando H2 revisado excluye el choque impuesto y los dos ticks de confirmación de la recuperación al medir el área de exceso después del retorno sostenido. Con profundidad 0.5, E3 menos E1-alta es 0.1229, 0.1239 y 0.1215 para los choques $-0.05$, $-0.10$ y $-0.20$; con profundidades 1 y 2, los contrastes pareados son cero a precisión numérica. La repetición de salidas explica por qué no bastaba una razón de flujo continuo: en E1 estático poco profundo, $R_{\mathrm{flow}}=0.670<1$, pero la razón del pánico único es $R_{\mathrm{panic}}=1.327>1$. Como medida conservadora de presión configurada en un tick, se usa $R_{\max}^{\mathrm{adv}}=1.586$, que incluye el máximo de ruido más pánico alineados; las razones componentes distinguen flujo persistente de un impulso único. Las salidas observadas de E1 coinciden con ese evento único. En E3, $R_{\mathrm{flow}}=1.208$ tras el retiro nominal y las salidas repetidas persisten sin pánico. Estas razones usan órdenes configuradas y alineadas, y reservas aproximadas; son diagnósticos explicativos, no reglas de umbral validadas. El control poco profundo con capacidad fija reduce el contraste del área posrecuperación a 0.0036--0.0097, lo cual muestra que el contraste original está sustancialmente ligado a la regla de capacidad basada en reservas en este diseño.

Para ingeniería, estos resultados justifican proponer campos de telemetría para futuros experimentos de reproducción: precios de referencia y piscina sincronizados, reservas activas, flujo de órdenes con signo, capacidad de arbitraje disponible, orden de eventos/transacciones, primera entrada a la banda, retorno sostenido y salidas posteriores. Es una propuesta de qué medir en el siguiente estudio, no evidencia de que una regla de monitoreo en vivo esté validada.

El aporte de la tesis es metodológico y acotado: deriva el referente de capacidad de producto constante, estudia la interacción de reglas seleccionadas en un simulador reproducible e informa sensibilidades y resultados negativos en la comparación de modelos. Una afirmación sobre sistemas reales exigiría reproducir eventos de mercado sincronizados, con comisiones del pool, gas, latencia, capital finito, orden transaccional y competencia entre buscadores. El manuscrito actual no incluye tal reproducción empírica.

\section{Limitaciones y amenazas a la validez}

El simulador usa un precio CEX exógeno, intercambios sin comisión, sin gas ni retraso de ejecución, inventarios amplios, un agente representativo por función, una piscina de producto constante y ticks de eventos sin duración de reloj calibrada. No representa posiciones de liquidez concentrada, competencia entre buscadores ni MEV. Las reglas de retiro y pánico son heurísticas fijas; el disparador H2 se activa en todas las corridas ensayadas y el retiro cambia tanto la profundidad como el límite de arbitraje. La sensibilidad y las comparaciones reservadas de modelos se mantienen dentro de esta única familia sintética de reglas. La comparación híbrida también usa la trayectoria futura realizada del CEX. No se incluye calibración, reproducción sincronizada de eventos ni validación externa.
